
\bigskip

\section{Types and the pro-completion}\label{pro}
In this section, we study the pro-completion of $\defT(T)$. In the category theory literature, Barr \cite{BARR1986113} studied $\Pro (\C)$ in the case when $\C$ is regular. Textbooks on shape theory also contain references for the pro-completion, such as \cites[Ch.~I \S1\addsemicolon]{alma992985533896801631}[Appendix]{alma992977584035101631}. 

\medskip

\subsection{Types and morphisms of types}
We first give the model theoretic description of $\Pro(\defT(T))$ as the category of types, which was first studied in \cite{KAMENSKY2007180}. Continue to fix $T$ to be a consistent complete first-order theory with monster model $\mathcal{U}$. Fix a suitably big cardinal $\kappa$ and say a set is small if its cardinality is strictly less than $\kappa$. When we say `sufficiently large', we mean larger than $\kappa$.

\begin{defin}
\begin{enumerate}
    \item A \emph{type} is a $T$-deductively-closed set of formulas in a small set of variables (not necessarily complete). To emphasise the categorical approach, we write $X, Y$ for a type (rather than the model-theoretic convention of $p,q$). We write $X_i$ for the formulas in the type $X$, and identify $(X_i)_{i\in I}$ with $X$ itself. The notation is justified by the fact that $X=\lim X_i$ in $\Pro(\defT(T))$.
    \item We define the \emph{category of types} as follows:
    \begin{itemize}
        \item Objects are types and falsity $\bot$;
        \item A morphism $f:(X_i)_{i\in I}\rightarrow (Y_j)_{j\in J}$ is induced by a collection of definable functions $f_j: X_{i_j} \rightarrow Y_j$ for each $j\in J$, which is compatible in the following sense: for any $j,j'\in J$,
    
        \[ T \cup X(x) \vdash f_j(x)=f_{j'}(x) \]

        We then take a quotient on the set of compatible collections of definable functions by the following equivalence relation: for any two compatible collections $(f_j)_{j\in J}, (g_j)_{j\in J}$ from $X$ to $Y$,

        \[(f_j)\sim (g_j) \ \Leftrightarrow \ T \cup X(x) \vdash f_j(x)=g_j(x) \textup{ for all }j. \]
        Hence, more precisely a morphism is an equivalence class of compatible collections of definable functions.
    \end{itemize}
\end{enumerate}
\end{defin}

\begin{remark}
    We require a type to be deductively closed, so that if two morphisms in the category of types are $T$-equivalent, then they are the same morphism. In particular, this ensures that there is only one identity morphism for each object, and composition is well-defined. Having said that, we might define a specific type by giving a non-deductively-closed set of formulas. Formally, we mean the deductive closure of the given set.
    
    Unpacking the definition of a morphism in the special case when the domain or the codomain is $\bot$, one sees that $\bot$ is the initial object.

    By compactness, a collection of definable functions is compatible if and only if for any $j,j'\in J$, there is some $k\in J$ such that $T\cup X_k(x)\vdash f_j(x) = f_{j'}(x)$. Likewise, two compatible collections $(f_j),(g_j)$ are equivalent if and only if for all $j$, there is some $k$ such that $T\cup X_k(x)\vdash f_j(x)=g_j(x)$. Compare this to \cite[Ch.~I, pp. 4--6]{alma992985533896801631}.

    A compactness argument also shows that a morphism is precisely a pro-definable set, whose interpretation in a suitably saturated $\U$ is the graph of a function. This is a result of \cite{KAMENSKY2007180}, which is recorded here as a lemma.
\end{remark}

\begin{lem}[{\cite[Corollary 9]{KAMENSKY2007180}}]\label{functions of Pro}
    Let $\U\models T$ be sufficiently saturated. Let $X,Y\in\Pro(\defT(T))$. There is a natural bijection between $\Hom(X,Y)$ and the subobjects of $X\times Y$, whose interpretation in $\mathcal{U}$ is a function from $X(\mathcal{U})$ to $Y(\mathcal{U})$
\end{lem}

\smallskip

We recall the categorical construction of $\Pro (\C)$.
\begin{defin}
\begin{enumerate}
    \item A small category $I$ is \emph{cofiltered} if it has finite cones.
    \item A \emph{cofiltered limit} in $\C$ is a limit whose diagram is given by a cofiltered category.
    \item Given a cofiltered diagram $F:I\rightarrow\C$, the \emph{transition maps} of $F$ are the morphisms $F\alpha:Fi\rightarrow Fj$ for morphisms $\alpha:i\rightarrow j$ in $I$.
    \item Given a category $\C$, the \emph{pro-completion} $\Pro(\C)$ is a category with cofiltered limits, equipped with a functor $\C\rightarrow\Pro(\C)$ inducing, for any $\mathcal{D}$ with cofiltered limits, an equivalence of categories
        \[ \Fun(\C,\mathcal{D}) \simeq \textup{Cofilt}(\Pro(\C),\mathcal{D}),\]
        where $\textup{Cofilt}$ denotes the full subcategory of the functor category spanned by cofiltered-limit-preserving functors.
\end{enumerate}
\end{defin}

The universal property defines $\Pro(\C)$ up to equivalence. When $\C$ has all finite limits, $\Pro(\mathcal{C})$ is equivalent to $\textup{Lex}(\mathcal{C},\textbf{Set})^{op}$, the opposite category of finite-limit-preserving functors into \textbf{Set}. The embedding $\C\rightarrow\Pro(\C)$ is the restricted Yoneda embedding. This follows from \cite[Definition~6.1.1, Remark~6.1.7]{lurie_ultracategories}; see also \cite[Appendix, Corollary 2.8]{alma992977584035101631}.

Alternatively, by \cite[Remark 6.1.6]{lurie_ultracategories}, $\Pro(\C)$ admits the following description:
\begin{itemize}
    \item Objects are cofiltered diagrams of $\C$;
    \item On morphisms, we have the natural bijection 
    \[ \hom_{\Pro(\mathcal{C})}((X_i)_{i\in I}, (Y_j)_{j\in J}) \cong \lim_{j\in J} \colim_{i\in I} \hom_\mathcal{C}(X_i,Y_j). \]
\end{itemize}

\begin{lem}
    The category $\Pro(\defT(T))$ is equivalent to the category of types.
\end{lem}

\begin{proof}
    On objects, given a cofiltered system $(X_i)_{i\in I}$, we can form the a type in the variables $(x_i)_{i\in I}$, where each $x_i$ is pairwise disjoint, by taking the deductive closure of

    \[ \{ X_i(x_i) \ | \ i\in I \} \cup \{ \alpha(x_i)=x_j \ | \ \alpha: X_i \rightarrow X_j \textup{ is a transition map}\}. \]
    Conversely given a type $\{\phi_i(x_i) \ | \ i\in I \}$ in the variables $x=\bigcup x_i$, we define the cofiltered system $X_i = \phi_i(x)$, where there is an inclusion morphism $X_i\rightarrow X_j$ whenever $T\vdash X_i\rightarrow X_j$.

    On morphisms, the equivalence is simply spelling out the natural bijection; for details, see \cite[Ch.~I~\S1.1]{alma992985533896801631}.
\end{proof}

\begin{cor}
    The category of types is a regular category and the embedding $\defT(T)\rightarrow\Pro(\defT(T))$ is a fully faithful regular embedding.
\end{cor}

\begin{proof}
    By \cite[Theorem 1]{BARR1986113}, this is true for any regular category, in particular $\defT(T)$.
\end{proof}

\smallskip

Historically, people studying pro-completions have been interested in \emph{uniform approximation}. The problem is as follows: given a certain diagram of shape $I$ in $\Pro(\C)$, can we find diagrams of the same shape in $\C$, whose cofiltered limit is the given diagram? More explicitly, consider that the category $[I,\Pro(\C)]$ have cofiltered limits and there is a functor $[I,\C]\rightarrow[I,\Pro(\C)]$ by composing with the inclusion functor $\C\rightarrow\Pro(\C)$. The universal property of $\Pro([I,\C])$ then induces a unique functor 
\[E^I_\C: \Pro([I,\C]) \rightarrow [I,\Pro(\C)].\] 
The uniform approximation problem then asks: under what assumptions is $E^I_\C$ an equivalence? This was studied in \cites[Expos\'{e} I Proposition 8.8.5 \addsemicolon]{a.TheorieToposCohomologie1972}[Appendix, Proposition 3.3 \addsemicolon]{alma992977584035101631}[Ch.~II \S3]{meyerphd}. The problem turns out to be quite subtle, and as the most recent contribution \cite[Example 1.5]{henryIndcompletionAccessibilityFunctor2025} points out, there have been mistakes in the past literature. We state the uniform approximation theorem, specialised and dualised to fit our setting.

\begin{thm}[Uniform approximation {\cite[Theorems 1.2 and 1.3]{henryIndcompletionAccessibilityFunctor2025}}]\label{simonhenry}
    Let $I$ be an indexing category.
    \begin{enumerate}
        \item The category $I$ is finite if and only if for any category $\C$ with finite limits, the functor $E^I_\C: \Pro([I,\C]) \rightarrow [I,\Pro(\C)]$ is an equivalence.
        \item The category $I$ is finite and has no non-identity endomorphism if and only if for any category $\C$, the functor $E^I_\C: \Pro([I,\C]) \rightarrow [I,\Pro(\C)]$ is an equivalence.
    \end{enumerate}
\end{thm}

As a corollary, we obtain an easier description of the morphisms of $\Pro(\C)$. Given a morphism $f$ in $\Pro(\C)$, the domain $X$ and the codomain $Y$ can be given by the same indexing category as $(X_i)$ and $(Y_i)$, respectively. Moreover, there are morphisms $(f_i:X_i\rightarrow Y_i)$ in $\C$, whose limits in the arrow category $[\{\bot\rightarrow \top\},\C]$ is $f$.

\begin{defin}
    A morphism $f:X\rightarrow Y$ in $\Pro(\mathcal{C})$ is called a \textit{level morphism} if there is a cofiltered category $I$ with a functor $F:I\rightarrow [\{\bot\rightarrow \top\},\C]$, such that $\lim_I F \cong f$. In this case, we write $f\cong\lim(f_i:X_i\rightarrow Y_i)$, where $f_i = Fi$ for $i\in I$.
\end{defin}

\begin{lem}\label{level morphisms}
    Every morphism in $\Pro(\mathcal{C})$ is a level morphism. 
\end{lem}

\begin{proof}
    Apply part (2) of Theorem \ref{simonhenry} to $I=\{\bot\rightarrow \top\}$.
\end{proof}

\smallskip

Let us now study when two objects of $\Pro(\C)$ are isomorphic. 

\begin{defin}[{\cite[p.~149]{alma992977584035101631}}]
Let $I$ and $J$ be cofiltered categories. We say a functor $F:I\rightarrow J$ is \emph{cofinal} if for every $j\in J$, there is $i\in I$ and a morphism $Fi\rightarrow j$ in $J$.
\end{defin}

\begin{lem}[{\cite[Appendix, Corollary 2.5]{alma992977584035101631}}]\label{cofinal functor}
    Let $F:I\rightarrow J$ be a cofinal functor between cofiltered categories. Then, the pro-objects $(X_j)_{j\in J}$ and $(X_{Fi})_{i\in I}$ are isomorphic via the representing morphisms $(\textup{id}_{X_{Fi}})_{i\in I}$.
\end{lem}

\begin{prop}[{\cite[Proposition~8]{KAMENSKY2007180}}]\label{isomorphism is calculated in monster}
    A morphism in $\Pro(\defT(T))$ is an isomorphism if and only if its interpretation in some/any saturated model is a bijection. In other words, saturated models reflect isomorphisms.
\end{prop} 

\begin{remark}
    Any model $\M: \defT(T)\rightarrow \textbf{Set}$ extends uniquely to a continuous functor $\M: \Pro(\defT(T))\rightarrow \textbf{Set}$. Hence, an isomorphism of $\Pro(\defT(T))$ is mapped to an isomorphism of \textbf{Set}, which is a bijection. The other direction requires a use of compactness on formulas with negation. In particular, the proposition does not hold if $T$ is merely a positive theory. 
\end{remark}


\begin{example}
    We show that non-saturated models might not reflect isomorphisms. Let $\mathcal{L} = \{c_n \ | \ n\in\omega\}$ be a countably infinite set of constants. Let $T$ be the theory which says each $c_n$ is distinct. Then, $p := \{ x \neq c_n \ | \ n\in\omega \}$ is a consistent type. Consider the empty map into $p$.

    Let $\M$ be the prime model, where every element is named by some $c_n$, so that $p$ is not realised. Then, the empty map is a bijection. It is not an isomorphism: using the proposition above, a saturated model realises $p$, so that the empty map is not a bijection.
\end{example}


The following lemma allows us to ignore the differences between categorical limits and set-theoretic intersections.

\begin{lem}\label{level intersections}
    Let $f:X\rightarrow Y$ be a morphism in $\Pro(\defT(T))$. Then, there are formulas $\phi_i,\psi_i$ and definable functions $g_i:\phi_i\rightarrow\psi_i$, such that in any model $\M\models T$, we have $X(\M)=\bigcap \phi_i(\M)$, $Y(\M)=\bigcap \psi_i(\M)$, and $f(\M)=\bigcap g_i(\M)$.
\end{lem}

\begin{proof}
    By Lemma \ref{level morphisms}, assume $f=\lim (f_i:X_i\rightarrow Y_i)$. Moreover, let $\phi_\alpha$ enumerate all the formulas which are consequences of $X$, each considered in the same set of variables as $X$. Likewise let $\psi_\beta$ enumerate the consequences of $Y$. By construction, $X=\bigcap \phi_\alpha$ and $Y=\bigcap \psi_\beta$ in any model.

    Given $i\in I$, the map $\pi_i: X\rightarrow X_i$ implies that $X_i(\pi_i(x))$ is one of the consequences of $X$, so is one of the $\phi_\alpha$. Conversely, given $\alpha$, compactness implies that there is some $i\in I$ such that $X_i(\pi_i(x))\vdash \phi_\alpha(x)$, so that there is a morphism $X_i\rightarrow\phi_\alpha$. This is true for $Y_i$ and $\psi_\beta$ as well. Hence, given $i\in I$, we form the following diagram, which commutes with the projections from $X$ and $Y$.
        % https://q.uiver.app/#q=WzAsNyxbMCw0LCJYX2kiXSxbMSw0LCJZX2kiXSxbMCwyLCJYX2oiXSxbMSwyLCJZX2oiXSxbMSwzLCJcXHBzaV97XFxiZXRhX2l9Il0sWzAsMSwiXFxwaGlfe1xcYWxwaGFfaX06PVhfaihcXHBpX2ooeCkpIl0sWzAsMF0sWzAsMSwiZl9pIiwyXSxbMiwzLCJmX2oiXSxbNCwxXSxbMiwwXSxbMyw0XSxbNSwyLCJcXHBpX2oiLDJdXQ==
        \[\begin{tikzcd}
        	{} & \\
        	{\phi_{\alpha_i}:=X_j(\pi_j(x))} \\
        	{X_j} & {Y_j} \\
        	& {\psi_{\beta_i}} \\
        	{X_i} & {Y_i}
        	\arrow["{\pi_j}"', from=2-1, to=3-1]
        	\arrow["{f_j}", from=3-1, to=3-2]
        	\arrow[from=3-1, to=5-1]
        	\arrow[from=3-2, to=4-2]
        	\arrow[from=4-2, to=5-2]
        	\arrow["{f_i}"', from=5-1, to=5-2]
        \end{tikzcd}\]
    Let $g_i$ be the morphism $\phi_{\alpha_i}\rightarrow\psi_{\beta_i}$. Then as $(g_i)$ is cofinal in $(f_i)$, they have the same limit. 
\end{proof}

\medskip

\subsection{Regular epimorphisms}
We already know that $\Pro(\defT(T))$ is regular. We study its regular epimorphisms and obtain a characterisation of saturated models. We first recall the \emph{diagonal fill-in property} of a regular category.

\begin{prop}[{\cite[Proposition 2.4]{Butz1998}}]\label{diagonal fill in}
    In a regular category $\C$, given a commutative square
    % https://q.uiver.app/#q=WzAsNCxbMCwwLCJYIl0sWzAsMSwiWCciXSxbMSwwLCJZIl0sWzEsMSwiWSciXSxbMCwxLCJlIiwyLHsic3R5bGUiOnsiaGVhZCI6eyJuYW1lIjoiZXBpIn19fV0sWzIsMywibSIsMCx7InN0eWxlIjp7InRhaWwiOnsibmFtZSI6Imhvb2siLCJzaWRlIjoidG9wIn19fV0sWzAsMiwiZiJdLFsxLDMsImciLDJdXQ==
    \[\begin{tikzcd}
    	X & Y \\
    	{X'} & {Y'}
    	\arrow["f", from=1-1, to=1-2]
    	\arrow["e"', two heads, from=1-1, to=2-1]
    	\arrow["m", hook, from=1-2, to=2-2]
    	\arrow["g"', from=2-1, to=2-2]
    \end{tikzcd}\]
    where $e$ is a regular epimorphism and $m$ is a monomorphism, there is a unique arrow $d:X'\rightarrow Y$ such that the two triangles commute.
\end{prop}

\begin{lem}\label{level regular epimorphism}
    Let $\mathcal{C}$ be regular. Then a morphism in $\Pro(\mathcal{C})$ is a regular epimorphism if and only if it is a cofiltered limit of regular epimorphisms in $\mathcal{C}$.
\end{lem}

\begin{proof}
    $(\Rightarrow)$ \quad The dual statement was proved in \cite[Theorem 1]{BARR1986113}. Let $f=\lim f_i$. Given $i$, we may factorise $f_i$ into a regular epimorphism $e_i$ followed by a monomorphism $m_i$. Denote by $W_i$ the image of this factorisation. By the diagonal fill-in property in Proposition~\ref{diagonal fill in}, we have a unique morphism $Y\rightarrow W_i$, which exhibits $Y$ as the limit of $W_i$.
    % https://q.uiver.app/#q=WzAsNSxbMCwwLCJYIl0sWzEsMCwiWF97aV97al8wfX0iXSxbMiwwLCJXIl0sWzIsMSwiWV97al8wfSJdLFswLDEsIlkiXSxbMCwxXSxbMSwyLCJlXzAiLDAseyJzdHlsZSI6eyJoZWFkIjp7Im5hbWUiOiJlcGkifX19XSxbMiwzLCJtIiwwLHsic3R5bGUiOnsidGFpbCI6eyJuYW1lIjoiaG9vayIsInNpZGUiOiJ0b3AifX19XSxbMCw0LCJmIiwyLHsic3R5bGUiOnsiaGVhZCI6eyJuYW1lIjoiZXBpIn19fV0sWzQsM10sWzQsMiwiIiwxLHsic3R5bGUiOnsiYm9keSI6eyJuYW1lIjoiZGFzaGVkIn19fV1d
        \[\begin{tikzcd}
            X & {X_i} & {W_i} \\
            Y && {Y_i}
            \arrow[from=1-1, to=1-2]
            \arrow["f"', two heads, from=1-1, to=2-1]
            \arrow["{e_i}", two heads, from=1-2, to=1-3]
            \arrow["{m_i}", hook, from=1-3, to=2-3]
            \arrow[dashed, from=2-1, to=1-3]
            \arrow[from=2-1, to=2-3]
        \end{tikzcd}\]
    Then, $f$ is the cofiltered limit of $e_i$. This shows that all regular epimorphisms in $\Pro(\C)$ are cofiltered limits of regular epimorphisms in $\C$.

\smallskip

    $(\Leftarrow)$ \quad  In \textbf{Set}, filtered colimits commute with finite limits. Since limits and colimits of presheaves are calculated pointwise, this holds also for $\textup{Lex}(\mathcal{C},\textbf{Set}) \simeq \Pro(\mathcal{C})^{op}$. Dualising, in $\Pro(\mathcal{C})$ cofiltered limits commute with finite colimits. In particular, cofiltered limits of regular epimorphisms are regular epimorphisms. 
\end{proof}

Warning: it is not true that if $(f_i)$ represents a regular epimorphism $f$, then each $f_i$ is a regular epimorphism. Instead, the lemma says that we can find some representing $(e_i)$, each of which is a regular epimorphism.

\smallskip

The above proof also gives a description of the regular epimorphism-monomorphism factorisation in $\Pro(\mathcal{C})$. Let $f=\lim f_i$. Then, one factorises each $f_i = m_i \circ e_i$ through its image $W_i$. Then, using the diagonal fill-in property, whenever we have $j\rightarrow i$, there is an induced $W_j \rightarrow W_i$. Hence, we can take the limit $\lim W_i$, which gives precisely the image of $f$ with the regular epimorphism $\lim e_i$ and the monomorphism $\lim m_i$. This was first shown in \cite{BARR1986113}. For a more modern account, see \cite[Theorem 12]{kanalas2025puremapsstrictmonomorphisms}.
% https://q.uiver.app/#q=WzAsOSxbMCwyLCJYX2kiXSxbMiwyLCJZX2kiXSxbMSwyLCJXX2kiXSxbMCwxLCJYX2oiXSxbMSwxLCJXX2oiXSxbMiwxLCJZX2oiXSxbMCwwLCJYIl0sWzEsMCwiXFxsaW0gV19pIl0sWzIsMCwiWSJdLFswLDIsImVfaSIsMix7InN0eWxlIjp7ImhlYWQiOnsibmFtZSI6ImVwaSJ9fX1dLFsyLDEsIm1faSIsMix7InN0eWxlIjp7InRhaWwiOnsibmFtZSI6Imhvb2siLCJzaWRlIjoidG9wIn19fV0sWzQsNSwibV9qIiwyLHsic3R5bGUiOnsidGFpbCI6eyJuYW1lIjoiaG9vayIsInNpZGUiOiJ0b3AifX19XSxbMyw0LCJlX2oiLDIseyJzdHlsZSI6eyJoZWFkIjp7Im5hbWUiOiJlcGkifX19XSxbMywwXSxbNSwxXSxbNCwyLCIiLDEseyJzdHlsZSI6eyJib2R5Ijp7Im5hbWUiOiJkYXNoZWQifX19XSxbOCw1XSxbNyw0XSxbNiwzXSxbNiw3LCJcXGxpbSBlX2kiLDIseyJzdHlsZSI6eyJib2R5Ijp7Im5hbWUiOiJkYXNoZWQifSwiaGVhZCI6eyJuYW1lIjoiZXBpIn19fV0sWzcsOCwiXFxsaW0gbV9pIiwyLHsic3R5bGUiOnsidGFpbCI6eyJuYW1lIjoiaG9vayIsInNpZGUiOiJ0b3AifSwiYm9keSI6eyJuYW1lIjoiZGFzaGVkIn19fV0sWzYsOCwiZiIsMCx7ImN1cnZlIjotMn1dXQ==
\[\begin{tikzcd}
	X & {\lim W_i} & Y \\
	{X_j} & {W_j} & {Y_j} \\
	{X_i} & {W_i} & {Y_i}
	\arrow["{\lim e_i}"', dashed, two heads, from=1-1, to=1-2]
	\arrow["f", curve={height=-12pt}, from=1-1, to=1-3]
	\arrow[from=1-1, to=2-1]
	\arrow["{\lim m_i}"', dashed, hook, from=1-2, to=1-3]
	\arrow[from=1-2, to=2-2]
	\arrow[from=1-3, to=2-3]
	\arrow["{e_j}"', two heads, from=2-1, to=2-2]
	\arrow[from=2-1, to=3-1]
	\arrow["{m_j}"', hook, from=2-2, to=2-3]
	\arrow[dashed, from=2-2, to=3-2]
	\arrow[from=2-3, to=3-3]
	\arrow["{e_i}"', two heads, from=3-1, to=3-2]
	\arrow["{m_i}"', hook, from=3-2, to=3-3]
\end{tikzcd}\]
With $\mathcal{C} = \defT(T)$, the image of a morphism $f:X\rightarrow Y$ is given by the type 
\[Y \cup \{\exists x \ (X_{i_j}(x) \wedge f_j(x)=y ) \ | \ j\in J\}.\]

\smallskip

We now recall the notion of $\lambda$-regularity.

\begin{defin}[{\cite[Definition 2.2]{MAKKAI1990225}}]
    Let $\lambda$ be a regular cardinal. 
\begin{enumerate}
    \item A category $\C$ is $\lambda$-\textit{regular} (resp. $\infty$-\emph{regular}) if it is regular, has $\lambda$-limits (resp. small limits), and regulars epimorphisms are preserved by $\lambda$-products (resp. small products).
    \item  A category $\C$ is $\lambda$-\textit{exact} (resp. $\infty$-\emph{exact}) if it is $\lambda$-regular (resp. $\infty$-regular) and is exact.
    \item  A functor between $\lambda$-regular categories is $\lambda$-\emph{regular} (resp. $\infty$-\emph{regular}) if it preserves regular epimorphisms and $\lambda$-limits (resp. small limits)
\end{enumerate}
\end{defin}

\begin{thm}\label{pro is infty regular}
    Let $\C$ be regular. Then $\Pro(\C)$ is $\infty$-regular.
\end{thm}

\begin{proof}
    We already know that $\Pro(\C)$ is regular and complete. Moreover, a cofiltered limit of regular epimorphisms is a regular epimorphism. Since $\C$ has finite limits, every small product in $\Pro(\C)$ is a cofiltered limit of representibles. Hence, a product of regular epimorphisms is a cofiltered limit of regular epimorphisms, which is itself a regular epimorphism.

    More explicitly, let $f_i: X_i \rightarrow Y_i$ be regular epimorphisms for each $i\in I$, where $I$ is an infinite but small set. Let $J$ be the poset category of non-empty finite subsets of $I$. Then $J^{op}$ is a cofiltered category. For each $j=\{i_1,\ldots,i_n\}\in J$, let $g_j = \prod_{k=1}^{n} f_{i_k}$. Then, $\lim_{j\in J^{op}} g_j = \prod_{i\in I} f_i$.
\end{proof}

\begin{remark}
    The Axiom of Choice is equivalent to the statement that, in \textbf{Set}, infinite products of regular epimorphisms are regular epimorphisms. Theorem~\ref{pro is infty regular} says that the pro-completion of any regular category satisfies the Axiom of Choice in this sense.
\end{remark}

\medskip

\subsection{Saturation}\label{subsect_saturation}

We take a brief detour to obtain a characterisation of saturation. We first give a motivating example.

\begin{example}
    Recall that an object $X$ is \textit{inhabited} if the unique morphism $X\rightarrow \top$ into the terminal object is a regular epimorphism. In $\defT(T)$, an object is not inhabited if and only if it is $\bot$. We prove this also holds in $\Pro(\defT(T))$. It is clear that the morphism $\bot\rightarrow \top$ is not a regular epimorphism in $\Pro(\defT(T))$. Conversely, suppose $f:X\rightarrow \top$ is not a regular epimorphism in $\Pro(\defT(T))$. Note that by Lemma~\ref{functions of Pro}, we have $X\cong\bot$ if and only if there is a morphism $X\rightarrow\bot$. Now, $f$ is represented by some $f_i:X_i\rightarrow \top$. If all the $f_i$ are regular epimorphisms, then $f$ is also a regular epimorphism by Lemma~\ref{level regular epimorphism}. Hence, there exists some $i_0$ such that $f_{i_0}$ is not a regular epimorphism, so in fact $X_{i_0} \cong \bot$. There is a projection map $X\rightarrow X_{i_0}$, so in fact $X\cong\bot$.

    In $\defT(T)$, an inhabited object is always non-empty in any model, but in $\Pro(\defT(T))$, an inhabited object can be empty in a model by the omitting type theorem. Categorically, a model $\M$ is a coherent functor $\defT(T)\rightarrow\textbf{Set}$; in particular $\M$ preserves regular epimorphisms in $\defT(T)$. Now, $\M$ admits a unique extension to a continuous functor $\Pro(\defT(T))\rightarrow \textbf{Set}$, but this extension might not be preserve regular epimorphisms in $\Pro(\defT(T))$. 
\end{example}

A $\kappa$-saturated model is a model which realises all types with sets of parameters strictly smaller than $\kappa$. Accordingly, we need a way of measuring the size of a set of parameters. 

\begin{defin}
    The \textit{length} of an object $X\in\Pro(\C)$, denoted by $l(X)$, is the least cardinality $|I|$ such that there is a monomorphism $X\hookrightarrow \prod_{i\in I}Y_i$, where $Y_i\in\C$.

    The length of a morphism in $\Pro(\C)$ is the length of its codomain.
\end{defin}

\begin{lem}
    Suppose $X\in\Pro(\defT(T))$ has infinite length. Then $X$ is isomorphic to a type $p(x)$, where $l(X)=|x|$.
\end{lem}

\begin{proof}
    Let $\lambda=l(X)$, so that there is a monomorphism $m:X\hookrightarrow\prod_{\alpha<\lambda} Y_\alpha$, where $Y_\alpha\in\defT(T)$ is a formula $\psi_\alpha(y_\alpha)$. Let $p(x)$ be the image of the regular epimorphism-monomorphism factorisation of $m$. Since $m$ is monomorphism, $X$ is isomorphic to $p(x)$. Moreover, the image is in the same tuple of variable as $\prod_{\alpha<\lambda} Y_\alpha$, which is $(y_\alpha)_{\alpha<\lambda}$. This has length $\lambda$. 
\end{proof}

\begin{remark}
    An object with finite length in fact has length 1. Not every type in a finite tuple of variables can be expressed as a type in a single variable. However, this is not a big problem: every type in a finite tuple of variables can be expressed as a type in a single variable in $T^{eq}$.
\end{remark}

\begin{remark}\label{cocompact}
    Another way of measuring the size of an object $X$ in a category (for example $\Pro(\C)$) is by considering whether the functor $\Hom(-,X)$ preserves $\lambda$-cofiltered limits. If so, then $X$ is called $\lambda$-\emph{cocompact}. The \emph{size} of $X$ is the least $\lambda$ such that $X$ is $\lambda$-cocompact. This roughly corresponds to the cardinality of the smallest cofiltered diagram landing in $\C$, whose limit in $\Pro(\C)$ is $X$. We say an object is \emph{cocompact} if it is $\aleph_0$-cocompact. By \cite[Theorem 7.2]{DanielC2002}, an object is cocompact if and only if it is isomorphic to an object in $\C$. We write $\textup{Cocompact}(\C)$ for the full subcategory of $\C$ consisting of the cocompact objects. Then, we can recover $\C$ from $\Pro(\C)$ uniquely via $\textup{Cocompact}(\Pro(\C)) \simeq \C$. Specialising to $\defT(T)$, an object is cocompact if and only if it is an isolated type.
    
    The size matches with the length when they are larger than $|T|$. However, when they are smaller, they can be different. For example, a non-isolated type in a finite tuple has length 1, but infinite size.
\end{remark}


\smallskip

We show that saturation is equivalent to preservation of regular epimorphisms in $\Pro(\defT(T))$.

\begin{thm}
    A model $\M$ is $\lambda$-saturated, if and only if the functor $\M:\Pro(\defT(T))\rightarrow \textbf{Set}$ maps regular epimorphisms with length strictly less than $\lambda$ to surjections. 
\end{thm}

\begin{proof}
    $(\Rightarrow)$ \quad  Let $f=\lim (f_i: X_i\rightarrow Y_i)$ be a morphism $X\rightarrow Y$ in $\Pro(\defT(T))$, where $l(Y)<\lambda$. By Lemma \ref{level regular epimorphism}, assume all $f_i$ are regular epimorphisms. We write $p_i$ for the projection $X\rightarrow X_i$ and $\pi_i$ for the projection $Y\rightarrow Y_i$.

    Let $b\in Y(\M)$ realise $Y$ in $\M$, with $|b|\leq l(Y)<\lambda$. Then, the type $X\cup\{f_i(p_i(x))=\pi_i(b) \ |\ i\in I \}$ is consistent, since each $f_j$ is a surjection. Since $\M$ is $\lambda$-saturated, the type is realised by some $a\in X(\M)$. It is clear that $f^\M(a)=b$.

\smallskip

    $(\Leftarrow)$ \quad Let $p(x)\in S(A)$ be a type over some parameter set $A\subset \M$, where $|A|<\lambda$. Let $X(x,y_A)$ be the type $p'\cup \tp(A)$, where $\tp(A)$ is in the variables $y_A$ and $p'$ is obtained from $p$ by replacing $A$ with variables $y_A$. Then, that $\M$ realises $p(x)$ will follow once we show the surjectivity of the projection $\pi_A: X\rightarrow \tp(A)$. Note that $l(\pi_A) = l(\tp(A)) \leq |A| < \lambda$. It suffices to show that $\pi_A$ is a cofiltered limit of regular epimorphisms.

    The map $\pi_A$ is a cofiltered limit of the maps $[\phi(x,y_A) \wedge \psi(y_A)] \rightarrow [\psi(y_A)]$ for each $\phi\in p'$ and $\psi\in\tp(A)$. By the completeness of $T$, we have
    \[T \vdash \forall y_A \ (\psi(y_A) \rightarrow \exists x \ \phi(x,y_A)).\]
    Hence, these maps are indeed regular epimorphisms.
\end{proof}


\begin{remark}
    A coherent functor from $\Pro(\defT(T))$ to \textbf{Set} is therefore a $\lambda$-saturated model for all small $\lambda$. Monster models are precisely the homogeneous coherent functors. 

    Another way of characterising saturated models categorically is via a lifting property. $\M$ is $\lambda$-saturated if and only if for all models $\N\rightarrow \N'$, where $|\N'|< \lambda$, and elementary embedding $\N\rightarrow \M$, there is a lift $\N'\rightarrow \M$.
    % https://q.uiver.app/#q=WzAsMyxbMCwwLCJOJyJdLFswLDEsIk4iXSxbMSwwLCJNIl0sWzEsMF0sWzEsMl0sWzAsMiwiIiwxLHsic3R5bGUiOnsiYm9keSI6eyJuYW1lIjoiZGFzaGVkIn19fV1d
    \[\begin{tikzcd}
        {\N'} & \M \\
        \N
        \arrow[dashed, from=1-1, to=1-2]
        \arrow[from=2-1, to=1-1]
        \arrow[from=2-1, to=1-2]
    \end{tikzcd}\]
\end{remark}

We give an example to illustrate why saturated models are relevant.

\begin{example}[Infinitary Chinese remainder theorem]
    Let $T$ be the theory of $(\mathbb{Z},0,1,+,\times)$. For every prime $p$, the equivalence relation $\textup{mod } p$, as well as its quotient $\mathbb{Z}_p$ and quotient map $\pi_p$ is definable. Let $\pi = \langle \pi_p \rangle_p :\M\rightarrow \prod_p \mathbb{Z}_p$. This map is a cofiltered limit of regular epimorphisms $\pi_p$, so is a regular epimorphism. However, the interpretation in the standard model is not a surjection; for example, the tuple $(1,1,\ldots)$ does not have a pre-image. In a non-standard model $\M$, the tuple does have a pre-image whenever $\M$ realises the type $\{ x = 1 \textup{ mod } p\ |\ p \textup{ prime}\}$.
\end{example}

